\paragraph{\textbf{Remark 3.9.}}
\SGARefTarget{sga3:xx:3.9}
We may summarize a number of the results of this section by the
following diagram of principal homogeneous bundles (on the left):

\begin{center}
  \begin{minipage}{\textwidth}
  \centering
  \begin{tikzcd}[column sep=huge]
    \mathrm W(\mathfrak g^\alpha)^\times
      \arrow[r,"w_\alpha"]
      & N^\times
        \arrow[r,"a_\alpha"]
      & \mathrm W\bigl((\mathfrak g^\alpha)^{\otimes 2}\bigr)^\times
  \end{tikzcd}

  \smallskip
  \begin{tikzcd}[column sep=huge]
    \mathbb G_{m,S}
      \arrow[r,"\alpha^\ast"]
      & T
        \arrow[r,"\alpha"]
      & \mathbb G_{m,S}
  \end{tikzcd}
  \end{minipage}
\end{center}

Notice that \(a_\alpha\) is faithfully flat (since \(\alpha\) is),
and that \(w_\alpha\) is a monomorphism if and only if
\(\alpha^\ast\) is a monomorphism.  We leave to the reader the task
of writing down the corresponding diagrams for right
principal-bundle structures, analogous diagrams for the root
\(-\alpha\), and of studying the relations among these various
diagrams.

\paragraph{\textbf{Lemma 3.10.}}
\SGARefTarget{sga3:xx:3.10}
\emph{Let \(S\) be a scheme, and let \(q>0\) be an integer such that
\(x\mapsto x^q\) defines an endomorphism of \(\mathbb G_{a,S}\).
Let \((G,T,\alpha)\) and \((G',T',\alpha')\) be two elementary
\(S\)-systems, and let \(f\colon G\to G'\) be a morphism of
\(S\)-groups.  Let}
\[
  h\colon(\mathfrak g^\alpha)^{\otimes q}
    \xrightarrow{\ \sim\ }\mathfrak g'^{\,\alpha'}
\]
\emph{be an isomorphism of \(\mathcal O_S\)-modules, and let}
\[
  h^\vee\colon(\mathfrak g^{-\alpha})^{\otimes q}
    \xrightarrow{\ \sim\ }\mathfrak g'^{-\alpha'}
\]
\emph{be the contragredient isomorphism.  Suppose that, for every
\(S'\to S\) and every \(X\in W(\mathfrak g^\alpha)(S')\),}
\[
  f\bigl(\exp(X)\bigr)=\exp\bigl(h(X^q)\bigr).
\]
\emph{Then the following conditions are equivalent:}

\smallskip
\noindent\emph{\textup{\SGAInlineRefTarget{sga3:xx:3.10:item:i}(i)}
\(f\bigl(\alpha^\ast(z)\bigr)=\alpha'^\ast(z)^q\).}

\smallskip
\noindent\emph{\textup{\SGAInlineRefTarget{sga3:xx:3.10:item:ii}(ii)}
\(f\bigl(w_\alpha(Z)\bigr)=w_{\alpha'}\bigl(h(Z^q)\bigr)\).}

\smallskip
\noindent\emph{\textup{\SGAInlineRefTarget{sga3:xx:3.10:item:iii}(iii)}
\(f\bigl(\exp(Y)\bigr)=\exp\bigl(h^\vee(Y^q)\bigr)\).}

\smallskip
Each condition is to be read as holding for every \(S'\to S\) and
every
\[
  z\in\mathbb G_m(S'),\qquad
  Z\in W(\mathfrak g^\alpha)^\times(S'),\qquad
  Y\in W(\mathfrak g^{-\alpha})(S'),
\]
respectively.  Indeed, \textup{\SGARefLink{sga3:xx:3.8:item:i}{(i)}} implies \textup{\SGARefLink{sga3:xx:3.8:item:ii}{(ii)}} by \SGARefLink{sga3:xx:3.8}{3.8},
\textup{\SGARefLink{sga3:xx:3.8:item:ii}{(ii)}} implies \textup{\SGARefLink{sga3:xx:3.10:item:iii}{(iii)}} by \SGARefLink{sga3:xx:3.7}{3.7}, and \textup{\SGARefLink{sga3:xx:3.10:item:iii}{(iii)}}
implies \textup{\SGARefLink{sga3:xx:3.8:item:i}{(i)}} by \SGARefLink{sga3:xx:2.7}{2.7}.

\paragraph{\textbf{Proposition 3.11.}}
\SGARefTarget{sga3:xx:3.11}
\emph{Let \(S\) be a scheme, and let \(q\in\mathbb Z\), \(q>0\), be
such that \(x\mapsto x^q\) defines an endomorphism of
\(\mathbb G_{a,S}\).  Let \((G,T,\alpha)\) and
\((G',T',\alpha')\) be two elementary \(S\)-systems, and let
\(f\colon G\to G'\) be a morphism of \(S\)-groups.  The following
conditions on \(f\) are equivalent:}

\smallskip
\noindent\emph{\textup{\SGARefLink{sga3:xx:3.10:item:i}{(i)}} The restriction of \(f\) to \(T\)
factors through a morphism \(f_T\colon T\to T'\) making the diagram}

\begin{center}
  \begin{tikzcd}[column sep=huge,row sep=large]
    \mathbb G_{m,S}
      \arrow[r,"\alpha^\ast"]
      \arrow[d,"q"']
      & T
        \arrow[r,"\alpha"]
        \arrow[d,"f_T"']
      & \mathbb G_{m,S}
        \arrow[d,"q"'] \\
    \mathbb G_{m,S}
      \arrow[r,"\alpha'^\ast"]
      & T'
        \arrow[r,"\alpha'"]
      & \mathbb G_{m,S}
  \end{tikzcd}
\end{center}

\noindent\emph{commutative.}

\smallskip
\noindent\emph{\textup{\SGAInlineRefTarget{sga3:xx:3.11:item:ii}(ii)} There exists a unique isomorphism of
\(\mathcal O_S\)-modules}
\[
  h\colon(\mathfrak g^\alpha)^{\otimes q}
    \xrightarrow{\ \sim\ }\mathfrak g'^{\,\alpha'}
\]
\emph{such that}
\[
  f\bigl(\exp(X)\bigr)=\exp\bigl(h(X^q)\bigr),
  \qquad
  f\bigl(\exp(Y)\bigr)=\exp\bigl(h^\vee(Y^q)\bigr)
\]
\emph{for every \(S'\to S\),
\(X\in W(\mathfrak g^\alpha)(S')\), and
\(Y\in W(\mathfrak g^{-\alpha})(S')\).  Consequently, \(f\) also
satisfies the equivalent conditions of \SGARefLink{sga3:xx:3.10}{3.10}.}

\smallskip
We first prove that \textup{\SGARefLink{sga3:xx:3.10:item:ii}{(ii)}} implies \textup{\SGARefLink{sga3:xx:3.10:item:i}{(i)}}.  By \SGARefLink{sga3:xx:3.10}{3.10},
condition \textup{\SGARefLink{sga3:xx:3.10:item:ii}{(ii)}} gives
\[
  f\circ\alpha^\ast=q\alpha'^\ast;
\]
hence, by \SGARefLink{sga3:xx:3.3}{3.3}, \(f|_T\) factors through \(T'\).  It remains to prove
\[
  \alpha'\bigl(f(t)\bigr)=\alpha(t)^q,
\]
which follows at once from the fact that \(f\) induces a morphism of
groups
\[
  T\cdot U_\alpha\longrightarrow T'\cdot U_{\alpha'}.
\]

We now prove that \textup{\SGARefLink{sga3:xx:3.10:item:i}{(i)}} implies \textup{\SGARefLink{sga3:xx:3.11:item:ii}{(ii)}}.  Let
\[
  X\in\Gamma(S,\mathfrak g^\alpha),
  \qquad
  Y\in\Gamma(S,\mathfrak g^{-\alpha}),
\]
and set
\[
  p_+(x)=f\bigl(\exp(xX)\bigr),
  \qquad
  p_-(y)=f\bigl(\exp(yY)\bigr).
\]
These are morphisms of groups
\[
  p_+,p_-\colon\mathbb G_{a,S}\longrightarrow G'.
\]
Now
\[
\begin{aligned}
  \intaut\bigl(\alpha'^\ast(z)\bigr)^q\bigl(p_+(x)\bigr)
    &={}\intaut\bigl(f_T(\alpha^\ast(z))\bigr)
       \bigl(f(\exp(xX))\bigr)\\
    &={}f\Bigl(\intaut\bigl(\alpha^\ast(z)\bigr)
       \bigl(\exp(xX)\bigr)\Bigr)\\
    &={}f\bigl(\exp(z^2xX)\bigr)
     =p_+(z^2x).
\end{aligned}
\]
Applying \SGARefLink{sga3:xx:1.2}{Lemma 1.2} with
\(Q=\alpha'^\ast(\mathbb G_{m,S})\), we obtain a section
\[
  X'\in\Gamma(S,\mathfrak g'^{\,\alpha'})
\]
such that
\[
  f\bigl(\exp(xX)\bigr)=p_+(x)=\exp(x^qX').
\]
Similarly, there exists a section
\[
  Y'\in\Gamma(S,\mathfrak g'^{-\alpha'})
\]
such that
\[
  f\bigl(\exp(yY)\bigr)=\exp(y^qY').
\]
Since \(f\) is a group morphism, it respects formula \textup{(F)};
hence
\[
  X^qY^q=(XY)^q=X'Y'.
\]
It follows readily that \(X^q\mapsto X'\) and \(Y^q\mapsto Y'\)
define the announced isomorphisms \(h\) and \(h^\vee\).
